% The baseline's command cycle: one turn of the Fuehrungskreislauf, the three
% ordered checks at the seam, and the execution model closing the loop.
% Source: workflow-baseline-README.md sec.4 and sec.4a,
% DOCTRINE_BASELINE_DESIGN.md sec.2.2 and sec.3.1. Input from
% inc/workflow_baseline.tex, Section "The Command Cycle and the Seam".
\begin{figure}[htbp]
\centering
\begin{tikzpicture}[
  stage/.style={draw, rounded corners=2pt, fill=blue!6, text width=5.2cm,
    minimum height=0.66cm, align=center, font=\footnotesize, inner sep=3pt},
  rung/.style={draw, rounded corners=2pt, fill=yellow!12, text width=5.2cm,
    minimum height=0.66cm, align=center, font=\footnotesize, inner sep=3pt},
  seam/.style={draw, thick, dashed, rounded corners=3pt, gray!60!black},
  ext/.style={draw, rounded corners=2pt, fill=green!6, align=center,
    font=\footnotesize, inner sep=4pt},
  db/.style={draw, thick, rounded corners=3pt, fill=orange!8, align=center,
    font=\footnotesize, inner sep=4pt},
  flow/.style={-{Stealth[length=1.8mm]}, thick},
  back/.style={-{Stealth[length=1.8mm]}, thick, dashed, red!55!black},
  lbl/.style={font=\scriptsize, align=center, text=gray!30!black},
]
  % --- the timeline and what feeds it --------------------------------
  \node[db, text width=5.2cm] (q) {\textbf{Timeline}\\ scenario events and replies,\\ in sim-time order};
  \node[ext, left=0.9cm of q, text width=1.7cm] (scen) {Scenario file};
  \draw[flow] (scen) -- (q);

  % --- Lagefeststellung ----------------------------------------------
  \node[stage, below=0.7cm of q] (lage) {\textbf{Lagefeststellung}\\ read the wake and\\ update the Lage (S2)};

  % --- Planung: the three ordered checks at the seam -----------------
  \node[rung, below=1.15cm of lage] (r1) {\textbf{1\enspace Contingency check (D8)}\\ does the committed plan still fit?};
  \node[rung, below=0.6cm of r1] (r2) {\textbf{2\enspace Fallback check (D9)}\\ is the Lage understood at all?};
  \node[rung, below=0.6cm of r2] (r3) {\textbf{3\enspace Phase ladder (D3)}\\ what is the next phase? (S3, S4)};
  \node[seam, fit=(r1)(r2)(r3), inner sep=6pt,
    label={[font=\footnotesize, anchor=south east, inner sep=2pt]north east:\textbf{Planung} (the seam)}] (plan) {};

  % --- Befehlsgebung and the world -----------------------------------
  \node[stage, below=0.9cm of plan] (bef) {\textbf{Befehlsgebung}\\ shared engine: check, commit,\\ send, write the event log};
  \node[ext, right=1.2cm of bef, text width=2.9cm] (world) {\textbf{Execution model}\\ and shared agents carry out the orders};

  % --- main path ------------------------------------------------------
  \draw[flow] (q) -- (lage);
  \draw[flow] (lage) -- (r1);
  \draw[flow] (r1) -- node[lbl, right] {plan fits} (r2);
  \draw[flow] (r2) -- node[lbl, right] {Lage understood} (r3);
  \draw[flow] (r3) -- node[lbl, right] {next phase, or hold} (bef);

  % --- early exits on the left ---------------------------------------
  \coordinate (xl) at ($(plan.west)+(-0.7,0)$);
  \draw[flow, rounded corners] (r1.west) -- (r1.west -| xl) |- (bef.west);
  \draw[thick] (r2.west) -- (r2.west -| xl);
  \node[lbl, left] at (r1.west -| xl) {mismatch:\\ corrective\\ order};
  \node[lbl, left] at (r2.west -| xl) {not understood:\\ holding\\ order};

  % --- orders out, replies back (Kontrolle) --------------------------
  \draw[flow] (bef) -- node[lbl, above] {orders} (world);
  \draw[back, rounded corners] (world.north) |- node[lbl, pos=0.3, right]
    {\textbf{Kontrolle}\\ replies join the\\ timeline when\\ they happen} (q.east);
\end{tikzpicture}
\caption[The baseline's command cycle]{One turn of the baseline's command cycle.
A wake is read into the situation picture (the \emph{Lage}). At the seam,
three deterministic checks run in a fixed order. A plan that no longer fits is
corrected first. A situation that is not understood is held next. Only then
does the phase ladder advance the mission. Everything from Befehlsgebung onward
is the agentic engine's own code. The execution model carries out the orders
and its replies re-enter the timeline, which closes the loop.}
\label{fig:wf:cycle}
\end{figure}
